\chapter{Аналитическая часть}
\section{Стандартный алгоритм умножения матриц}
Пусть:
\begin{equation}
	\label{eq:an_matrixies}
	A[n, m] = \begin{pmatrix}
		a_{11} & a_{12} & \ldots & a_{1n}\\
		a_{21} & a_{22} & \ldots & a_{2n}\\
		\vdots & \vdots & \ddots & \vdots\\
		a_{m1} & a_{m2} & \ldots & a_{mn}
	\end{pmatrix},
\end{equation}
\begin{equation}
	B[m, k] = \begin{pmatrix}
		b_{11} & b_{12} & \ldots & b_{1n}\\
		b_{21} & b_{22} & \ldots & b_{2n}\\
		\vdots & \vdots & \ddots & \vdots\\
		b_{k1} & b_{k2} & \ldots & b_{kn}
	\end{pmatrix}.
\end{equation}

Тогда результатом операции ($A \cdot B$) является матрица $C$, расчитанная по формуле:

\centerline{$c_{ij} = \sum_{r=1}^{m} A[i, r] B[r, j]$, где $(i=\overline{1,l}; j=\overline{1,n})$}

\section{Стандартный алгоритм Винограда}
В 1987 году Дон Копперсмит и Шмуэль Виноград опубликовали метод, содержащий асимптотическую сложность $O(n^{2,3755})$ и улучшили  его в 2011 до $O(n^{2,373})$, где $n$ -- размер сторон матрицы \cite{book_vinograd}.

Идея алгоритма заключается в снижении доли дорогих операций умножения.

Вводятся векторы $V = (v_1, v_2, v_3, v_4)$ и $W = (w_1, w_2, w_3, w_4)$, их скалярное произведение равно:

\begin{equation}
V \cdot W = v_1w_1 + v_2w_2 + v_3w_3 + v_4w_4.
\end{equation}
	
Это равенство можно переписать в виде:
\begin{equation}
	\label{eq:vecMul}
	V \cdot W = (v_1 + w_2)(v_2 + w_1) + (v_3 + w_4)(v_4 + w_3) - v_1v_2 - v_3v_4 - w_1w_2 - w_3w_4.	
\end{equation}

Пусть даны матрицы $A[n, m], B[m, k], C[n, k]$. Скалярное произведение (\ref{eq:vecMul}) можно использовать при расчете произведения матриц следующим образом:
\begin{multline} 
	\label{eq:vin_mul_matrix}
	C[i, j] = \sum_{k=1}^{q/2}(A[i, 2k-1] + B[2k, j])(A[i, 2k] + B[2k-1, j]) - \\
	- \sum_{k=1}^{q/2} A[i, 2k-1]A[i, 2k] - \sum_{k=1}^{q/2} B[2k-1, j]B[2k, j]
\end{multline}

В выражении (\ref{eq:vecMul}) операнды со знаком <<$-$>> могут быть вычислены заранее для каждой строки и каждого столбца исходных матриц и переиспользовны без необходимости их повторного вычисления.

Аналгично с (\ref{eq:vin_mul_matrix}), вычислив один раз $\sum_{k=1}^{q/2}A[i, 2k-1]A[i, 2k]$ для строки, данное значение можно использовать для вычисления результата любого произведения с операндами этой строки. При предварительном расчете значений столбцов используется аналогичная идея.

Если размерность матрицы нечётна, необходимо произвести дополнительные операции для учета последней строки и столбца.

\section{Оптимизированный алгоритм Винограда}

Возможны следующие оптимизации представленного выше алгоритма.
\begin{enumerate}
\item Использование операции двоичного сдвига вместо умножения на 2. 
\item Вынесение части рассчетов, производящихся в случае, если размерность матрицы нечётна, в общий цикл. Это позволяет снизить количество выполняемых операций за счет отсутствия необходимости использования ещё одного цикла.
\item Введение декремента при вычислении вспомогательных массивов. 
\end{enumerate}